\section{The remaining construction theorem}
\label{fmb:remaining}

The progress now has a precise measure-theoretic meaning. For one
specified completion of the local coarse map, the full finite
positive Wilson pushforward is the sum of positive history measures.
At the bare level, large-curvature cells and explicit fixed-box
gauge/minimizer defects have uniform-volume joint probability
bounds. Neither assertion discards a bad sector or a vacuum factor.

What remains is to make the two constructions refer to the same
iterated local coordinates and to control their interaction. A
useful next theorem is a \emph{complete-density transfer theorem}
for the actual chosen coarse map. It would need to provide:
\begin{enumerate}
\item A spatial good/bad decomposition on determining sets with
collars, prescribed exterior data, and retained holonomy/transition
data. Its good-component integral must equal the existing
constrained stationary-plus-fluctuation integral, with its actual
independent windows, Haar factors, pivot factors and vacuum
normalization.
\item A normalized replacement on each complementary component.
The numerator and denominator must be bounded together, uniformly
in the admissible exterior field, component geometry, preceding
depth and scale history. A bound sufficient for summation over
components and their histories is required. Bare probability
rarity or a Gaussian conditional estimate does not by itself
provide this normalized replacement.
\item Preservation of a specified complete-density class under the
full step, with local analytic/source domains and coupling
extraction belonging to that same class. The weak local shooting
from Notebook 37 can be used only after its recursion has been
identified with this full-density recursion, or after a shooting
theorem has been proved for the latter.
\end{enumerate}
The published large-field induction discussed above is a possible
route to this theorem after its hypotheses have been matched.
The new direct probability estimates are another input to a proof;
they do not already verify the published induction class.

\paragraph{Why local rarity is not a global chart theorem.}
For any fixed family of \(M\) boxes, the one-box defect estimate
gives a union bound proportional to \(M\). If the number of boxes
diverges with the cutoff, a small one-box probability does not
imply that all boxes are good. For example, on the local tuned
trajectories \(g_{\rm bare}\) has order \(n^{-1/2}\), while a
logarithmic window yields a displayed estimate of order
\(\exp[-c(\log n)^{2a}]\). This bound does not defeat a number
of microscopic boxes growing exponentially in \(n\). This is
a limitation of that union-bound inference, not a lower bound
on the actual defect probability. The purpose of a component
expansion is to allow defects and sum them with their correct
weights, rather than assume a globally good field.

\paragraph{After the complete-density theorem.}
Continuum existence still requires common estimates for a
sufficient gauge-invariant observable algebra, with renormalized
short-distance insertions and the exact same denominator. Those
estimates must produce a common subsequential continuum hierarchy,
reflection positivity inherited from one positive regulator
sequence, full Euclidean covariance, and the regularity needed
for reconstruction. A quantitative lower witness and
identification of the resulting fields and dynamics are needed
to establish a nontrivial Yang--Mills theory. The physical mass
gap is a subsequent assertion. The Gaussian gaps used in the
local analysis are not that spectral statement.

No uniqueness of the continuum limit is imposed merely for the
sake of an existence proof. Nor is a particular contact-term
convention required when an admissible consistent choice would
suffice. The unresolved requirements above concern the actual
measure, observable bounds and physical identification.

\section*{Verification record and source access}

The accompanying research notes preserve the exact finite-measure
construction, the self-contained Wilson chessboard proof, the
fixed-box conditional proof and their independent internal audits.
The assembled source was checked for duplicate and missing labels
and compiled with the desktop LaTeX compiler. These checks establish
document consistency and typesetting; the mathematical evidence is
the arguments and their stated hypotheses.

Primary references used to locate the complete-density interface are
T.~Ba{\l}aban, \emph{Convergent renormalization expansions for lattice
gauge theories}, Communications in Mathematical Physics
\textbf{119} (1988), 243--285,
\href{https://doi.org/10.1007/BF01217741}{doi:10.1007/BF01217741},
and \emph{Large field renormalization. II. Localization,
exponentiation, and bounds for the \(R\) operation},
Communications in Mathematical Physics \textbf{122} (1989),
355--392,
\href{https://doi.org/10.1007/BF01238433}{doi:10.1007/BF01238433}.
The exact Section 2 induction class in the first paper, and the
weak-coupling hypothesis in Theorem 1 of the second, are retained.
No unchecked OCR constants or powers are supplied as theorem
hypotheses. The first large-field paper's detailed replacement
formula has not been newly verified from an accessible original
in this pass; its exact audit remains part of any import route.

\end{document}
